% Basado en el template realizado por Diego Essaya, disponible en
%                                                         http://lug.fi.uba.ar
% Modificado por Sebastián Santisi.
% 2007: Modificado por Patricio Moreno y Michel Peterson.
% 2014: Modificado por Patricio Moreno.
% 2017: Modificado por Patricio Moreno.
% 2021: Modificado por Carla Sobico

% Acá se define el tamaño de letra principal:
% Para utilizar los estilos de KOMA-script, desconectar la línea siguiente y
% comentar la que le sigue (dejar sin comentar un único documentclass)
%\documentclass[10pt, spanish]{scrartcl}
\documentclass[a4paper, twoside, 10pt, spanish]{article}

%%%%%%%%%%%%%%%%%%%%%%%%%%%%%%
% CS
%%%%%%%%%%%%%%%%%%%%%%%%%%%%%%
\usepackage{booktabs}
\usepackage[margin=1in]{geometry}
\usepackage{array}
\usepackage{pdflscape}
% CONFIGURACIONES GENERALES
%%%%%%%%%%%%%%%%%%%%%%%%%%%%%%%%%%%%%%%%%%%%%%%%%%%%%%%%%%%%%%%%%%%%%%%%%%%%%
% Definición del tamaño de página y los márgenes:
% Si preferís menos márgenes, descomentá la línea siguiente
%\usepackage[a4paper,headheight=16pt,scale={0.7,0.8},hoffset=0.5cm]{geometry}
\usepackage[demo]{graphicx}
\usepackage{caption}
\usepackage{subcaption}
\usepackage{babel}  % contiene la correcta separación en sílabas del español
\usepackage[utf8x]{inputenc}    % porque el encoding del documento es UTF-8
\usepackage[per-mode=fraction]{siunitx}
\sisetup%
{
	output-decimal-marker = {,},
	exponent-product = \cdot,
    group-digits = integer,
	group-separator = {.}
}


%
% El paquete amsmath agrega algunas funcionalidades extra a las fórmulas.
% Además defino la numeración de las tablas y figuras al estilo "Figura 2.3",
% en lugar de "Figura 7". (Por lo tanto, aunque no uses fórmulas, si querés
% este tipo de numeración dejá el paquete amsmath descomentado).
%
\usepackage{amsmath, amsfonts, amssymb}
%\numberwithin{equation}{section}
%\numberwithin{figure}{section}
\numberwithin{table}{section}
%%%%%%%%%%%%%%%%%%%%%%%%%%%%%%%%%%%%%%%%%%%%%%%%%%%%%%%%%%%%%%%%%%%%%%%%%%%%%

%%%%%%%%%%%%%%%%%%%%%%%%%%%%%%%%%%%%%%%%%%%%%%%%%%%%%%%%%%%%%%%%%%%%%%%%%%%%%
% ENCABEZADO y PIE DE PÁGINA
%%%%%%%%%%%%%%%%%%%%%%%%%%%%%%%%%%%%%%%%%%%%%%%%%%%%%%%%%%%%%%%%%%%%%%%%%%%%%
\usepackage{fancyhdr}   % Para poder personalizarlo
\usepackage{lastpage}   % Para poder saber cuántas páginas tiene el documento
\pagestyle{fancy}
\renewcommand{\sectionmark}[1]{\markboth{}{\thesection\ \ #1}}
\fancyhead{}	% Elimino el contenido del encabezado
% Muestra la sección a la derecha (izquierda) en páginas impares (pares)
\fancyhead[RO,LE]{\rightmark}
% El siguiente texto a la derecha (izquierda) en páginas pares (impares)
\fancyhead[RE,LO]{TB068 - Entrega \Nro~1}
\fancyfoot{}	% Elimino el contenido del pie de página
% A la izquierda (derecha) en páginas pares (impares): nro. de página / total
\fancyfoot[LE,RO]{\thepage/\pageref{LastPage}}
%%%%%%%%%%%%%%%%%%%%%%%%%%%%%%%%%%%%%%%%%%%%%%%%%%%%%%%%%%%%%%%%%%%%%%%%%%%%%

%%%%%%%%%%%%%%%%%%%%%%%%%%%%%%%%%%%%%%%%%%%%%%%%%%%%%%%%%%%%%%%%%%%%%%%%%%%%%
% Hipervínculos (enlaces) en el documento (y modificación de atributos)
%%%%%%%%%%%%%%%%%%%%%%%%%%%%%%%%%%%%%%%%%%%%%%%%%%%%%%%%%%%%%%%%%%%%%%%%%%%%%
\usepackage{url}
\urlstyle{tt}
\usepackage[colorlinks=true,linkcolor=black, urlcolor=blue]{hyperref}
\hypersetup{
    breaklinks,
    baseurl       = http://,
    pdfborder     = 0 0 0,
    pdfpagemode   = UseNone,
    pdfstartpage  = 1,
    pdfcreator    = {Plantilla de informe de TP para \LaTeX{}},
    bookmarksopen = true,
    bookmarksdepth= 2,% to show sections and subsections
    pdfauthor     = {Apellido~1, Apellido~2, Apellido~3},
    pdftitle      = {-},
    pdfsubject    = {Informe},
    pdfkeywords   = {}%
    }
%%%%%%%%%%%%%%%%%%%%%%%%%%%%%%%%%%%%%%%%%%%%%%%%%%%%%%%%%%%%%%%%%%%%%%%%%%%%%

%%%%%%%%%%%%%%%%%%%%%%%%%%%%%%%%%%%%%%%%%%%%%%%%%%%%%%%%%%%%%%%%%%%%%%%%%%%%%
% LISTAS (para poder modificar los 'bullets' de las listas)
%%%%%%%%%%%%%%%%%%%%%%%%%%%%%%%%%%%%%%%%%%%%%%%%%%%%%%%%%%%%%%%%%%%%%%%%%%%%%
\usepackage{enumerate}
%%%%%%%%%%%%%%%%%%%%%%%%%%%%%%%%%%%%%%%%%%%%%%%%%%%%%%%%%%%%%%%%%%%%%%%%%%%%%

%%%%%%%%%%%%%%%%%%%%%%%%%%%%%%%%%%%%%%%%%%%%%%%%%%%%%%%%%%%%%%%%%%%%%%%%%%%%%
% TABLAS (para que se vean bien)
%%%%%%%%%%%%%%%%%%%%%%%%%%%%%%%%%%%%%%%%%%%%%%%%%%%%%%%%%%%%%%%%%%%%%%%%%%%%%
\usepackage{booktabs}
\usepackage{multirow}
%%%%%%%%%%%%%%%%%%%%%%%%%%%%%%%%%%%%%%%%%%%%%%%%%%%%%%%%%%%%%%%%%%%%%%%%%%%%%
\usepackage{comment}
%%%%%%%%%%%%%%%%%%%%%%%%%%%%%%%%%%%%%%%%%%%%%%%%%%%%%%%%%%%%%%%%%%%%%%%%%%%%%
% IMÁGENES
%%%%%%%%%%%%%%%%%%%%%%%%%%%%%%%%%%%%%%%%%%%%%%%%%%%%%%%%%%%%%%%%%%%%%%%%%%%%%
% Para incluir imágenes, el siguiente código carga el paquete graphicx
% según se esté generando un archivo dvi o un pdf (con pdflatex).

% Para generar dvi, descomentá la linea siguiente:
%\usepackage[dvips]{graphicx}

% Para generar pdf, descomentá las dos lineas seguientes:
\usepackage{xcolor}
\definecolor{kicadpcbbg}{RGB}{0,16,35}
\usepackage{graphicx}
\pdfcompresslevel=9

% Todas las imágenes están en el directorio imgs:
\newcommand{\imgdir}{img}
\graphicspath{{\imgdir/}}
%%%%%%%%%%%%%%%%%%%%%%%%%%%%%%%%%%%%%%%%%%%%%%%%%%%%%%%%%%%%%%%%%%%%%%%%%%%%%

%%%%%%%%%%%%%%%%%%%%%%%%%%%%%%%%%%%%%%%%%%%%%%%%%%%%%%%%%%%%%%%%%%%%%%%%%%%%%
% DIAGRAMAS DE FLUJO EN DIA
%%%%%%%%%%%%%%%%%%%%%%%%%%%%%%%%%%%%%%%%%%%%%%%%%%%%%%%%%%%%%%%%%%%%%%%%%%%%%
% Necesitas este paquete si haces los diagramas de flujo en el programa Dia
% y exportás a latex
%\usepackage{tikz}
%%%%%%%%%%%%%%%%%%%%%%%%%%%%%%%%%%%%%%%%%%%%%%%%%%%%%%%%%%%%%%%%%%%%%%%%%%%%%

%%%%%%%%%%%%%%%%%%%%%%%%%%%%%%%%%%%%%%%%%%%%%%%%%%%%%%%%%%%%%%%%%%%%%%%%%%%%%
% INSERCIÓN DE CÓDIGO FUENTE
%%%%%%%%%%%%%%%%%%%%%%%%%%%%%%%%%%%%%%%%%%%%%%%%%%%%%%%%%%%%%%%%%%%%%%%%%%%%%
% El paquete recomendado actualmente es minted.
% Documentación: https://www.ctan.org/pkg/minted
\usepackage[
        section,    % Numera el código según la sección
    ]{minted}
% minted provee los comandos:
% 1)  \mint[<opciones>]{<lenguaje>}<delimitador><código><delimitador>
% 2)  \mintinline[<opciones>]{<lenguaje>}<delimitador><código><delimitador>
% 3)  \inputminted[<opciones>]{<lenguaje>}{<archivo>}
\setminted[c]{
%        style=,
        linenos,            % Mostrar los números de línea
        numberfirstline,    % Numerar SIEMPRE la primera línea mostrada
        tabsize=4,          % Reemplazar las tabulaciones por 4 espacios
        autogobble          % Eliminar espacio sobrante al comienzo
    }
%%%%%%%%%%%%%%%%%%%%%%%%%%%%%%%%%%%%%%%%%%%%%%%%%%%%%%%%%%%%%%%%%%%%%%%%%%%%%
% COMANDOS UTILES
%%%%%%%%%%%%%%%%%%%%%%%%%%%%%%%%%%%%%%%%%%%%%%%%%%%%%%%%%%%%%%%%%%%%%%%%%%%%%
% los siguientes comandos permiten escribir de manera uniforme en todo el
% documento

% Para poder manejar los espacios al final de los comandos propios
\usepackage{xspace}

% Abreviatura de 'número' utilizando letras voladas (correcto español)
\newcommand{\Nro}{N.\textsuperscript{o}\xspace}
\newcommand{\nro}{n.\textsuperscript{o}\xspace}
%%%%%%%%%%%%%%%%%%%%%%%%%%%%%%%%%%%%%%%%%%%%%%%%%%%%%%%%%%%%%%%%%%%%%%%%%%%%%
% PAQUETES EXTRAS
%%%%%%%%%%%%%%%%%%%%%%%%%%%%%%%%%%%%%%%%%%%%%%%%%%%%%%%%%%%%%%%%%%%%%%%%%%%%%
\usepackage{circuitikz}
\usepackage{float}
\usepackage{multicol}
%\usepackage{pdfpages}
%\usepackage{subfigure}
%\usepackage{graphicx}%
\usepackage{gensymb}


\usepackage{verbatim}
% \usepackage{listings}
% \usepackage{c/style} % include custom style for C.

% \definecolor{codegreen}{rgb}{0,0.6,0}
% \definecolor{codegray}{rgb}{0.5,0.5,0.5}
% \definecolor{codepurple}{rgb}{0.58,0,0.82}
% \definecolor{backcolour}{rgb}{0.95,0.95,0.92}
%
% \lstdefinestyle{mystyle}{
%     backgroundcolor=\color{backcolour},   
%     commentstyle=\color{codegreen},
%     keywordstyle=\color{magenta},
%     numberstyle=\tiny\color{codegray},
%     stringstyle=\color{codepurple},
%     basicstyle=\ttfamily\footnotesize,
%     breakatwhitespace=false,         
%     breaklines=true,                 
%     captionpos=b,                    
%     keepspaces=true,                 
%     numbers=left,                    
%     numbersep=5pt,                  
%     showspaces=false,                
%     showstringspaces=false,
%     showtabs=false,                  
%     tabsize=2,
%     aboveskip=0pt,
%     belowskip=0pt
% }
%
% \lstset{style=mystyle}


%%%%%%%%%%%%%%%%%%%%%%%%%%%%%%%%%%%%%%%%%%%%%%%%%%%%%%%%%%%%%%%%%%%%%%%%%%%%%
%%%%%%%%%%%%%%%%%%%%%%%%%%%%%%%%%%%%%%%%%%%%%%%%%%%%%%%%%%%%%%%%%%%%%%%%%%%%%
% INICIO DEL DOCUMENTO
%%%%%%%%%%%%%%%%%%%%%%%%%%%%%%%%%%%%%%%%%%%%%%%%%%%%%%%%%%%%%%%%%%%%%%%%%%%%%
%%%%%%%%%%%%%%%%%%%%%%%%%%%%%%%%%%%%%%%%%%%%%%%%%%%%%%%%%%%%%%%%%%%%%%%%%%%%%

\addto\captionsspanish{\renewcommand{\tablename}{Tabla}}
\begin{document}


%
% Carátula:
%
\begin{titlepage}

\thispagestyle{empty}

\begin{center}

\includegraphics[scale=0.3]{fiuba.pdf}\\
\large{\textsc{Universidad de Buenos Aires}}\\
\large{\textsc{Facultad de Ingeniería}}\\
% Modificar año y cuatrimestre
\small{Año 2025 - 2\textsuperscript{do} cuatrimestre}
\end{center}

\vfill

\begin{center} % Modificar el código de ser necesario
\Large{\underline{\textsc{Taller de Sistemas Embebidos}}}\\
\vspace{.5cm}
 \Large{\textsc{Trabajo Practico Final}}
\end{center}

\vfill

\begin{center} % Modificar el código de ser necesario
\Large{\underline{\textsc{Sistema de detección de apneas y variaciones de oxigenación}}}\\
\Large{\underline{\textsc{durante el sueño}}}\\
\end{center}

\vfill

 \begin{tabbing}


	FECHA : 4/03/2026\\% \today\\
 \\

 \\
	ESTUDIANTES: \hspace{-1cm}\=\+\hspace{1cm}\=\hspace{6cm}\=\\
 		Musso Carboni Tomás	\>\> 111548\\ 
              \>\footnotesize{\verb!tmusso@fi.uba.ar!}\\
        Alan	\>\> \\ 
              \>\footnotesize{\verb!fi.uba.ar!}\\
        Archetti Santiago \>\> 110707\\ 
              \>\footnotesize{\verb!sarchetti@fi.uba.ar!}\\
           
 \end{tabbing}



\vfill

\hrule
\vspace{0.2cm}

% Modificar código de ser necesario
\noindent\small{Taller de Sistemas Embebidos\hfill}

\end{titlepage}

%
% Hago que las páginas se comiencen a contar a partir de aquí:
%
\setcounter{page}{1}

%
% Pongo el índice en una página aparte:
%

\tableofcontents
\newpage
%
% Inicio del TP:
%

\section{Introducción general}
\subsection{Sistema de detección de apneas y variaciones de oxigenación durante el Sueño}
El sistema portátil de monitoreo respiratorio y de saturación de oxígeno desarrollado en este trabajo cuenta con:
\begin{itemize}
    \item Una placa con microcontrolador \textit{STM32 F103RB}, encargada de la adquisición, procesamiento digital de señales y control general del sistema.
    \item Un sensor óptico \textit{MAX30102}, que permite medir señales PPG mediante leds rojo e infrarrojo para el cálculo de $SpO_{2}$ y la estimación de la frecuencia respiratoria.
    \item Un módulo de comunicación Bluetooth \textit{HM-10}, que posibilita la transmisión inalámbrica de los datos hacia un dispositivo móvil para su visualización en tiempo real.
    \item Un display, que muestra datos por pantalla en tiempo real.
    \item Un sistema de alarmas locales compuesto por un buzzer y leds indicadores, que alertan al usuario ante condiciones críticas de saturación de oxígeno.
    \item Procesamiento digital de la señal implementado en lenguaje C.
    \item Lógica de control que activa alarmas cuando la saturación de oxígeno cae por debajo de los umbrales configurados.
    \item Modos de funcionamiento configurables (modo adulto, modo niño).
\begin{figure} [H]
    \centering
    \includegraphics[width=0.6 \linewidth]{diagrama_de_bloques.png}
    \caption{Diagrama de bloques del sistema}
    \label{fig:diag_bloques}
\end{figure}
\end{itemize}

\subsection{Análisis de sistemas similares al desarrollado}
La inspiración para el desarrollo de este proyecto se basa en el funcionamiento de un oxímetro de pulso, el cual utiliza señales PPG obtenidas mediante la emisión de luz roja e infrarroja para estimar la saturación de oxígeno en sangre.
Tomando como referencia este principio de funcionamiento, el presente proyecto propone extender sus capacidades, con el objetivo de desarrollar un sistema de monitoreo completo durante el sueño. Además de estimar la saturación de oxígeno,permite obtener y supervisar otros parámetros fisiológicos relevantes, como la frecuencia cardíaca, la frecuencia respiratoria y la detección de episodios de apnea, proporcionando así una herramienta más completa para el seguimiento del estado fisiológico del usuario durante el descanso.
\begin{figure}
    \centering
    \includegraphics[width=0.4\linewidth]{oximetro.png}
    \caption{Oxímetro de pulso}
    \label{fig:oximetro}
\end{figure}


%\begin{figure}
  %  \centering
   % 
\includegraphics[width=0.5\linewidth]{fiuba.pdf}
    %\caption{Enter Caption}
    %\label{fig:placeholder}
%\end{figure}
\section{Introducción específica}
\subsection{Descripción del equipo}
El dispositivo dispone de dos botones, uno de los cuales permite seleccionar el modo de funcionamiento.
Esta configuración es necesaria debido a que los valores normales de saturación de oxígeno en sangre, frecuencia respratoria y frecuencia cardíaca difieren según la edad del usuario.
Con esta consideración, se cuentan 2 modos de funcionamiento: adulto o pediátrico. 
Una vez establecidos los parametros correspondientes al modo seleccionado, el sensor óptico MAX30102 se coloca en el dedo del usuario.

\begin{figure}[h]
    \centering
    \includegraphics[height=0.3 \textwidth]{ppg}
    \hspace{10pt}
    \includegraphics[height=0.3 \textwidth]{ppg2}
    \caption{Funcionamiento del sensor MAX30102.}
    \label{fig:ppg2}
\end{figure}

Este sensor se utiliza para la medición de señales de fotopletismografía (PPG), el mismo integra dos dfuentes de luz LED (una roja  y una infrarroja), un fotodiodo y un sistema de conversión analógico-digita.
Debido a que el volumen de sangre en los vasos varía con cada latido cardíaco, la cantidad de luz absorbida también cambia de forma periódica.
Las variaciones de luz absorbida son captadas por el fotodiodo del sensor, el cual convierte la intensidad de la luz reflejada en una señal eléctrica.
Los valores digitalizados corresponden a las mediciones de luz detectada por los LEDs rojo e infrarrojo, los cuales constituyen la señal PPG.
Estos datos son transmitidos al microcontrolador mediante $I^2C$.
Posteriormente, las señales adquiridas son sometidas a una etapa de procesamiento digital, en la cual se aplican los algoritmos necesarios para calcular los valores de saturación de oxígeno en sangre, frecuencia cardíaca y frecuencia respiratoria.
% A continuación se detalla brevemente dicho procesamiento: 

Para la estimación de la frecuencia cardíaca, se utiliza principalmente la señal PPG infrarroja.
Se aplica un filtrado pasa-banda en el rango aproximado de 0.5 a 5Hz, con el objetivo de eliminar ruidos para posteriormente detectar los máximos locales de la señal filtrada, que corresponden a latidos válidos.
A partir del intervalo temporal entre picos consecutivos, se calcula la frecuencia cardíaca en latidos por minuto.
En el caso de la saturación de oxígeno en sangre ($SpO_{2}$), se utilizan ambas señales PPG. Cada una de estas señales se descompone en dos componentes, una componente continua (DC) y una componente alterna (AC). A partir de estas componentes se calcula una relación entre las señales roja e infrerroja.

\begin{equation}
    R = \frac{\frac{AC_{RED}}{DC_{RED}}}{\frac{AC_{IR}}{DC_{IR}}}
\end{equation}

\begin{equation}
    SpO_{2} \approx 110 -25*R
\end{equation}

Por último, la frecuencia respiratoria puede estimarse analizando las oscilaciones de baja frecuencia presentes en la señal infrarroja.
Se aplica un filtro pasa-bajos con una frecuencia cercana a 0.4 Hz.
Finalmente, mediante la detección de picos o el análisis de la frecuencia dominante de la señal filtrada, se obtiene el número de ciclos respiratorios por minuto. 
Luego de la etapa de procesamiento, y considerando los rangos previamente establecidos según el modo de funcionamiento seleccionado, el sistema transmite en tiempo real la información obtenida tanto al display como a un dispositivo movil en caso de haberse establecido la conexión mediante bluetooth.
Los datos transmitidos incluyen : 

\begin{itemize}
    \item Detección de apnea (presencia o ausencia). 
    \item Frecuencia respiratoria actual. 
    \item Frecuencia cardíaca actual. 
    \item Saturacióin de oxígeno en sangre actual.
\end{itemize}

En caso de que la frecuencia respiratoria descienda por debajo de los límites considerados peligrosos, el sistema activa el buzzer y un led para alertar al usuario.
Esta alarma puede desactivarse manualmente mediante el botón de apagado correpondiente.
Asimismo, el dispositivo incorpora un modo de bajo consumo energético, el cual se activa automáticamente cuando no se detecta contacto del usuario con el sensor durante un período prolongado de tiempo.

\subsection{Requerimientos funcionales}

\begin{table}[h]
    \centering
    \hspace*{-.8cm}
    % se podría quitar la columna ID
    \begin{tabular}{|c|c|c|c|}
        \hline
        Grupo & ID & Requerimiento Funcional & Descripción \\ 
        \hline
        Adquisición& 1.1 & Adquisición de señales & Capturar señales ópticas ROJO/IR del MAX30102 via I2C.\\ \hline
         & 1.2 &Lectura periódica& Obtener muestras desde la FIFO del sensor.\\ \hline 
         &1.3 & Análisis de variación& Detectar variaciones innecesarias para $SpO_{2}$ y respiración \\ \hline
        Procesamiento& 2.1 & Filtrado digital& Aplicar filtrado \\ \hline
        & 2.2 &Cálculo de $SpO_{2}$ & Calcular $SpO_{2}$ mediante el ratio R\\ \hline
        & 2.3 & Cálculo de respiración & Estimar frecuencia respiratoria mediante\\ \hline
        & 2.3 & Cálculo de respiración & análisis de envolvente o línea base. \\ \hline
         & 2.4 & Empaquetado de datos & Generar paquete con $SpO_{2}$ , respiración y estado. \\ \hline
        Indicadores& 3.1 & Alarma crítica & Activar alarma sonora cuando se detecta apnea. \\ \hline
         & 3.2 & Indicadores LED & Mostrar estados (crítico, normal, error) mediante LEDs. \\ \hline
        Comunicación & 4.1 & Bluetooth & Transmitir datos vía HM-10 por UART. \\ \hline
        & 4.2 & Envío de parámetros & Enviar $SpO_{2}$, respiraciones y estado del paciente a la app. \\ \hline
         & 4.3 & Registro de eventos & Informar alarmas, desconexiones o estados especiales. \\ \hline
        & 4.4 & Display & Mostrar parámetros de interes por pantalla.\\ \hline
         % & 4.5 & Recepción de & Recibir umbrales y parámetros desde la aplicación. \\
         % &  & configuración & Recibir umbrales y parámetros desde la aplicación. \\ \hline
        Alimentación& 5.1 & Alimentación & Alimentación vía USB \\ \hline
         & 5.2 & Bajo consumo & Entrar en modo Sleep cuando no hay actividad. \\ \hline
        & 5.3 & Reactivación & Retomar actividad ante movimiento o comando. \\ \hline
    \end{tabular}
    \caption{Tabla de requerimientos}
    \label{tab:tab_requerimientos}
\end{table}

\subsection{Casos de uso}
\subsection{Descripción de módulos externos utilizados}

\begin{itemize}
    \item Display LCD 16x2
    \item Sensor MAX30102
    \item Módulo Bluetooth HM-10
\end{itemize}

\subsubsection{Display LCD 16x2}
El display utilizado es un display LCD de 16x2 como se puede observar en la figura \ref{fig:LCD}.
Dichas conexiones se realizaron con una configuración de 4 bits.
El mismo efectúa la comunicación con la placa mediante pines \textit{GPIO}. 

\begin{figure} [H]
    \centering
    \includegraphics[width=0.5\linewidth]{lcd_display.png}
    \caption{Módulo 1602A}
    \label{fig:LCD}
\end{figure}

\subsubsection{HM-10}
Para la transmisión de los datos hacia el dispositivo móvil del usuario se emplea un módulo Bluetooth HM-10.
La comunicación entre el módulo Bluetooth y la placa de desarrollo se realiza usando la interfaz UART.
A través de esta interfaz serie, el microcontrolador transmite al módulo los parámetros de interés para su posterior envío al dispositivo móvil del usuario.

\begin{figure}[H]
    \centering
    \includegraphics[width=0.5\linewidth]{HM10.png}
    \caption{Módulo bluetooth HM-10}
    \label{fig:HM10}
\end{figure}

\subsubsection{MAX30102}
En cuanto a la conexión del módulo Bluetooth con la placa de desarrollo, se realizaron cuatro conexiones principales: alimentación (VCC), referencia a tierra (GND) y dos líneas de comunicación correspondientes a la interfaz serie.
Las líneas de datos se conectan mediante una configuración cruzada entre los pines de transmisión y recepción.
En particular, el pin TX (transmit) del módulo se conecta al pin RX (receive) del microcontrolador, mientras que el pin RX del módulo se conecta al pin TX del microcontrolador permitiendo de esta forma el intercambio de datos. 

\begin{figure}[H]
    \centering
    \includegraphics[width=0.5\linewidth]{MAX30102.png}
    \caption{Sensor MAX30102}
    \label{fig: MAX30102}
\end{figure}

\section{Diseño e Implementación}
\subsection{Hardware del sistema}
A continuación se muestra un diagrama de bloques del Hardware del equipo y se explicarán las conexiones para cada uno de los componentes. 
\subsubsection{Display LCD 16x2}
\begin{table}[H]
\centering
\begin{tabular}{|c|c|}
\hline
Nucleo & Display\\ \hline
5V& VSS \\ \hline
5V & VDD \\ \hline
B12 & RS \\ \hline
GND & RW \\ \hline
A11 & E \\ \hline
- & D0 \\ \hline
- & D1 \\ \hline
- & D2 \\ \hline
- & D3 \\ \hline
PB5 & D4 \\ \hline
PB4 & D5 \\ \hline
PB13 & D6 \\ \hline
PB14 & D7 \\ \hline
GND & K \\ \hline
\end{tabular}
\caption{Conexiones del Display}
\end{table}
El pin A se conecta a 5V a través de una resistencia de $1k\ohm$ mientras que el pin Vo se conecta a 5V y GND a través de un potenciometro de $100k\ohm$.
\subsubsection{Sensor MAX30102}
\begin{table}[H]
\centering
\begin{tabular}{|c|c|}
\hline
Nucleo & Sensor MAX30102 \\ \hline
GND & GND \\ \hline
3.3V & VCC \\ \hline
PB6 & SCL \\ \hline
PB7 & SDA \\ \hline
\end{tabular}
\caption{Conexiones del sensor MAX30102}
\end{table}

\subsubsection{Módulo bluetooth HM-10}
\begin{table}[H]
\centering
\begin{tabular}{|c|c|}
\hline
Nucleo & HM-10 \\ \hline
GND & GND \\ \hline
5V & VCC \\ \hline
PA3 & TXD \\ \hline
PA2 & RXD \\ \hline
\end{tabular}
\caption{Conexiones del módulo bluetooth HM-10}
\end{table}

\subsubsection{Buzzer}
\begin{table}[H]
\centering
\begin{tabular}{|c|c|}
\hline
Nucleo & Buzzer \\ \hline
GND & - \\ \hline
PC7 & + \\ \hline

\end{tabular}
\caption{Conexiones del sensor MAX30102}
\end{table}

\subsubsection{Botones}
\begin{table}[H]
\centering
\begin{tabular}{|c|c|}
\hline
Nucleo & BOTON \\ \hline
PC10 & SWITCH-MODE \\ \hline
PC11 & ALARM-OFF \\ \hline
\end{tabular}
\caption{Conexiones de los botones}
\end{table}
Siendo SWITCH-MODE y ALARM-OFF los botones encargados de cambiar el modo de uso y apagar la alarma respectivamente, ambos configurados en modo pull-up y conectados a GND del otro extremo.

\subsubsection{LEDs}
Para polarizar los LEDs con una fuente de 3.3V, las resistencias limitadoras se calculan mediante la siguiente fórmula: 
\begin{equation}
    R = \frac{V_{CC}- V_f}{I_f}
\end{equation}

\begin{itemize}
    \item $V_{CC}$ la tensión de alimentación (en este caso 3.3V).
    \item $I_f$ la corriente deseada por el LED(normalmente se utiliza entre 5ma y 10ma).
    \item $V_f$ la tensión directa del LED.
\end{itemize}
Dado que se utilizan LEDs de diferentes colores el valor de $V_f$ dependerá de la caída de tensión típica según el color: 
\begin{itemize}
    \item LED rojo $\longrightarrow$ indica estado de peligro junto con la alarma. 
    \item LED azul $\longrightarrow$ indica que la conexión bluetooth a sido exitosa.
    \item LED verde $\longrightarrow$ modo adulto. 
    \item LED amarillo $\longrightarrow$ modo niño.
\end{itemize}

\begin{table}[H]
\centering
\begin{tabular}{|c|c|}
\hline
Color LED & $V_f$ típico \\ \hline
Rojo & 2V \\ \hline
Amarillo & 2.1V \\ \hline
Verde & 2.2V \\ \hline
Azul & 3.1V \\ \hline
\end{tabular}
\caption{Caída de tensión típica según el color}
\end{table}
Se tomó como valor deseado de corriente $I_f = 10mA$.
Los valores de las resistencias utilizadas para los LEDs amarillo, verde y rojo son de $R = 220 \ohm$ y $R= 30\ohm$ para el LED azul, conectados como se observa en la figura \ref{fig:LEDs}. 

\begin{figure}[H]
    \centering
    \begin{circuitikz}
        \draw (-1,1) node{3,3V} (-.5,1) node[circ]{} -- (0,1)
        to [led] (2,1)
        to [R] (4,1) -- (4.5,1)
        -- ++(0,-1) node[tlground]{};
    \end{circuitikz}
    \caption{Conexión de los LEDs.}
    \label{fig:LEDs}
\end{figure}


\begin{table}[H]
    \centering
    \begin{tabular}{|c|c|} \hline
        Nucleo & LED \\ \hline
        PB10 & LED-Alarm \\ \hline
        PB11 & LED-Bluetooth \\ \hline
        PB2 & LED-ADULT \\ \hline
        PB1 & LED-KID \\ \hline
    \end{tabular}
    \caption{Conexión de los LEDs}
\end{table}

\subsection{Costo de componentes}

\begin{table}[H]
    \centering
    \begin{tabular}{|c|c|} \hline
        Componente & Costo\\ \hline
        MAX30102& \\ \hline
        HM10&  \\ \hline
        Display&  \\ \hline
        Resistencias& \\ \hline
        LEDs& \\ \hline
        Buzzer&  \\ \hline
        Botones&  \\ \hline
    \end{tabular}
    \caption{Costos de componentes}
\end{table}

\subsection{Firmware}
%
% El firmware del dispositivo se programó en lenguaje C.
% La división de responsabilidades se realizó en los siguientes módulos.

\subsubsection{Sensor}\label{sec:sensor} % (fold)
Este módulo representa un sensor que se comunica con un oxímetro de pulso (\textit{MAX30102}) y procesa la señal para obtener los datos en la siguiente estructura de datos.%~\ref{data_struct}.

%formatear mas lindo con colorcitos y eso
\begin{verbatim}
    typedef struct {
        uint8_t heart_rate;
        uint8_t respiratory_rate;
        uint8_t spo2;
        uint8_t apnea;
        uint32_t timestamp;
    } task_sensor_results_dta_t;
\end{verbatim}

La comunicación se realiza via $I^2C$ y en cada iteración se acceden los datos al buffer FIFO del MAX30102 para reconstruir la señal.
Una vez procesada, cada \texttt{DEL\_SEN\_MAX\_TICKS} se accede a la interfaz del sistema y se actualiza la estructura de datos, además de activar una variable booleana indicando que el valor no ha sido leído.

% subsubsection Sensor (end)

\subsubsection{Actuador}\label{sec:actuador} % (fold)
Consiste en los indicadores al usuario.
Consta de 4 estados.
Pueden estar prendidos, apagados, parpadeando o emtiendo un pulso.
Éstos últimos permiten ajustar la frecuencia con la que un led parpadea, o la frecuencia de un buzzer.
% subsubsection Actuador (end)

\subsubsection{Botones}\label{sec:botones} % (fold)
Modela el comportamiento de un pulsador mecanico.
Permite filrar por software el ruido durante los cambios de estado del botón (\textit{debouncing}).
Aunque el botón es mecánico, se puede modelar otras entradas digitales con ruido de la misma manera.
% subsubsection Botones (end)


\subsubsection{Sistema}\label{sec:sistema} % (fold)
Es la tarea principal encargada de recibir las entradas desde los botones y sensores y producir los efectos correspondientes en los otros módulos.
Como es relevante no perder ninguna entrada, cuenta con una interfaz que gestiona una cola de eventos.
Ésta, además, es donde el sensor guarda los datos que toma cada cierta cantidad de ticks.
En cada iteración, el sistema pregunta a la cola si hay algun evento disponible y si hay algun dato para enviar.
En caso de haber un dato disponible del sensor, lo envía por UART al HM10, que se encarga de transmitirlo por bluetooth, y se muestra por el display

% subsubsection Sistema (end)

\subsection{Diseño}\label{sub:diseño} % (fold)

Se utilizó \textit{KiCAD} durante la etapa de diseño, para preparar la pcb del dispositivo.

\begin{figure}[H]
    \centering
    \colorbox{kicadpcbbg}{\includegraphics[width=.55 \textwidth]{tp3.png}}
    \caption{Diseño de la PCB en KiCAD.}
\end{figure}


La disposición de la placa fue pensada para soldarse en una placa experimental.
Se representaron los cables como pistas y se tuvo en cuenta el espaciado entre orificios.

\begin{figure}[H]
    \centering
    \includegraphics[height=.8 \textwidth]{render_side.png}
    \hspace{4pt}
    \includegraphics[height=.8 \textwidth]{render.png}
    \caption{Previsualización 3D de la placa.}
    \label{fig:}
\end{figure}

Se cuenta con espacio suficiente para agregar una batería para hacer al dispositivo portable.
Considerando que solo se precisa el integrado y el uso de una pantalla con menor tamaño y similar o mejor resolución queda margen para un modelo más compacto. 

{\color{red}en conclusiones tal vez?}

% subsection Diseño (end)

\subsection{Ensamblaje}\label{sub:ensamblaje} % (fold)
% subsection Ensamblaje (end)

\section{Ensayos y resultados}\label{sec:ensayos_y_resultados} % (fold)

% section Ensayos y resultados (end)


\end{document}
